\section{Consistency checks}

A useful skeleton should degrade sensibly when one of its mechanisms is removed.

\subsection{No competence evolution}

If
\[
\mathcal L(S_t,x_t,\theta_t)=S_t,
\]
operation does not change competence. The feedback through learning disappears and the
problem reduces to repeated organization of a fixed knowledge/competence state. This is the
natural static limit of Beam~1.

\subsection{No effect of allocation on experience}

If realized experience \(x_t\) is independent of \(a_t\), then current work organization does
not influence future competence through the proposed interface. The two beams coexist but are
not dynamically coupled. This is an important null case.

\subsection{No depreciation}

Setting retention to one in a representation such as
\[
e_{t+1}(q)=\rho_q e_t(q)+x_t(q)
\]
removes forgetting while retaining learning by experience. This corresponds to a standard
limiting case of competence-development models.

\subsection{No learning}

Setting the learning contribution to zero eliminates the Beam~2 feedback. A dynamic controller
should then have no competence-development reason to sacrifice current operational value.

\subsection{Why these checks matter}

These reductions show that the combined skeleton does not require learning to manufacture an
effect. If the relevant dynamic mechanisms are absent, the dynamic coupling disappears. This
is a basic but necessary safeguard against an artificial construction.
